\bookchapter{제28장: 돈을 알맞게 써서 마음에 드는 기기 사기——스마트폰 구매 궁극 가이드}

\chapterintro{스마트폰을 사는 일에 대해, 많은 사람들이 일주일 전에 리뷰를 보고, 한 달 가격을 비교하지만, 결과적으로 여전히 헷갈리며 5-6천 위안을 써서 사실 자신에게는 필요 없는 플래그십을 삽니다. 이 챕터는 파라미터 표를 나열하지 않습니다——그런 것은 공식 웹사이트에서 모두 볼 수 있음. 제가 하고 싶은 것은 “기기 선택 사고방식”을 드리는 것이며, 이후 시장에 어떤 새 스마트폰이 나와도, 그것이 자신에게 적합한지 스스로 판단할 수 있도록.}

\booksection{세 단계 기기 선택법: 요구예산기종}

많은 사람들이 인터넷에서 “3천 위안에 어떤 스마트폰 살까”라고 묻지만, 이 질문은 사실 답할 수 없습니다——왜냐하면 당신이 스마트폰을 무엇에 주로 사용하는지 모르기 때문. 스마트폰 구매의 올바른 사고 순서는 다음과 같아야 합니다:

\begin{tcbinfo}{기기 선택 세 가지 도끼}
\begin{enumerate}
  \item \textbf{먼저 가장 중시하는 것 명확히}: 촬영? 게임? 배터리 수명? 화면? 시스템 원활함? 아니면 전화·위챗만 가능하면 됨? 요구를 순서대로 정렬.
  \item \textbf{그 다음 예산 상한 정하기}: "대략 얼마"가 아니라, "최대 얼마나 쓸 수 있는지". 단단한 상한을 정해야 기기를 보면서 점점 비싼 것으로 가지를 않음.
  \item \textbf{마지막에 기종 선택}: 요구 순서와 예산을 가지고 기종을 매칭.
\end{enumerate}

\end{tcbinfo}

\booksection{핵심 파라미터 어떻게 보는지——모두 알 필요 없이, 이 5가지만 알면 충분}

파라미터 표에 수십 가지 데이터가 있지만, 일반 사용자가 주의해야 할 것은 이 5가지 차원뿐:

\begin{center}
\begin{tabularx}{\linewidth}{|X|X|X|}
\hline
\textbf{차원} & \textbf{어떻게 보는지} & \textbf{한 문장 요약} \tabularnewline
\hline
프로세서(SoC) & 스마트폰이 몇 년간 원활하게 사용할 수 있는지 결정. 2026년 기준: 스냅드래곤 8 Elite Gen5/디멘시티 9500 이상=최고급; 스냅드래곤 8s Gen4/디멘시티 8400=중고급; 스냅드래곤 7+ Gen4/디멘시티 7350=3년 사용 가능. 스냅드래곤 6 시리즈 미만은 신중 고려. & 새 것이 옳다, SoC가 스마트폰이 몇 년간 안 끊어지는지 결정 \tabularnewline
화면 & OLED는 LCD보다 색감이 좋고 더 절전; 주사율 120Hz는 60Hz보다 화면 스크롤이 훨씬 매끄러움; 해상도 최소 1080P 이상. 고주파 PWM 디밍(1920Hz 이상)은 눈에 더 친화적. & OLED+120Hz+고주파 디밍=좋은 화면의 합격선 \tabularnewline
카메라 & \textbf{픽셀 수가 아닌 센서 크기를 볼 것}. 센서가 더 크면 압도적, 센서 크기가 픽셀보다 훨씬 중요. 초광각은 실용적; 망원(3x 이상 광학 줌)은 멀리 있는 것을 촬영; 접사 렌즈는 대부분 메꾸기. & 대형 센서 메인 카메라+실용 망원>여러 메꾸기 렌즈 \tabularnewline
배터리 및 충전 & 4500mAh는 기본 합격선, 5000mAh 이상은 하루 한 번 충전 매우 안정. 급속 충전은 와트로: 60W 이상은 빠름, 100W 이상은 매우 빠름. 무선 충전은 편리하지만 느림. & 5000mAh+60W 급속 충전=불안 없는 조합 \tabularnewline
시스템 업데이트 약속 & 스마트폰을 얼마나 오래 안전하게 사용할 수 있는지 결정. Apple은 일반적으로 5-6년 시스템 업데이트; 삼성 플래그십은 4년 시스템+5년 보안 업데이트; Google Pixel은 7년; 국내 중급기 업데이트 약속은 일반적으로 짧음(2-3년). & 오래 사용하려면, 시스템 업데이트 약속이 긴 브랜드 선택 \tabularnewline
\hline
\end{tabularx}
\end{center}

\booksection{예산별 등급 추천(2026년 중 시장 기준)}

다음 추천은 2026년 중국 본토 시장 시세를 기준으로 하며, 가격은 공식 출시가. 구체적인 가격은 변동될 수 있으며, 실제 구매 시점의 가격을 기준으로 함.

\booksubsection{1000위안 이하: 전화 가능, 위챗 사용 가능, 충분}

\textbf{대상}: 학생, 보조 기기, 스마트폰에 대한 요구 높지 않음, 주로 전화와 위챗 사용.

\textbf{살 수 있는 것}: Redmi Note 시리즈, Honor X 시리즈, realme 숫자 시리즈의 중저가 모델. 프로세서는 대략 스냅드래곤 6 시리즈 또는 디멘시티 6000 시리즈, 화면은 LCD 위주, 촬영은 찍을 수 있지만 포커스가 느림, 시스템 원활도 보통.

\textbf{구매 가이드}: 이 가격대에서 두 가지 볼 것: 배터리가 충분히 큰지(5000mAh 이상 선택), 메모리가 충분한지(6GB+128GB부터 시작, 4GB는 끊기므로 구매 금지). Redmi가 이 가격대에서 가장 유리.

\booksubsection{1000-2000위안: 가성비 최고, 3-4년 안 끊어지는 스윗 스팟}

\textbf{대상}: 학생, 예산 제한적이나 오래 사용하고 싶음, 촬영과 게임에 극한 추구 없음.

\textbf{대표 기종}: Redmi Turbo 시리즈, iQOO Z 시리즈, Honor X 시리즈 고급, OnePlus Ace 시리즈 입문.

\textbf{구매 가이드}: 스냅드래곤 7+ Gen3 또는 디멘시티 7300 이상 프로세서를 탑재한 기종 선택, 8GB+256GB 메모리 구성. 이 구성은 일상 사용+가끔 왕자영광(왕저롱야오) 플레이에 충분, 정상적으로 3-4년 안 끊김.

\booksubsection{2000-3500위안: 중견 세력, 성능과 촬영 모두 뚜렷한 향상}

\textbf{대상}: 괜찮은 촬영을 원하고, 게임이 안 끊기길 바라며, 오래 사용하고 싶은 일반 사용자. 가장 많은 사람이 구매하는 가격대.

\textbf{대표 기종}: 샤오미 숫자 시리즈(예: 샤오미 17 표준판 가격 인하 후), vivo S 시리즈 고급, OPPO Reno 시리즈, Honor 숫자 시리즈, Redmi K 시리즈.

\textbf{구매 가이드}: 이 가격대의 가장 큰 특징은 “물통 기기”——뚜렷한 단점 없음. 2,000위안 이하 기기와 비교하여, 주요 향상은: 카메라가 뚜렷하게 좋아짐(광학 손떨림 보정 추가), 화면에 고급형 OLED, 급속 충전 더 빠름(80W 이상 흔함), 공예 품질도 더 정교. 12GB+256GB 또는 12GB+512GB 선택, 4-5년 사용 문제 없음.

\booksubsection{3500-5000위안: 준 플래그십, 게이커 외에는 모두 충분}

\textbf{대상}: 촬영에 추구가 있는 사용자, 최고 품질을 체험하고 싶은 사용자.

\textbf{대표 기종}: 샤오미 숫자 시리즈 표준판, vivo X 시리즈, OPPO Find X 시리즈 표준판, iPhone 구형 플래그십(예: iPhone 16).

\textbf{구매 가이드}: 이 가격대에서 이미 작년 플래그십 기종을 살 수 있음. 제 조언은:\textbf{올해 중급이 아닌 작년 플래그십을 사라}. 작년 플래그십의 재료, 공예, 카메라 하드웨어는 일반적으로 올해 중급보다 한 등급 높음. 전제는 “최신”이 아니어도 받아들일 수 있다는 것.

\booksubsection{5000위안 이상: 플래그십 구간, 돈에 상관없거나 특정 경험에 극한 추구}

\textbf{대상}: 사진 애호가, 중증 게임 플레이어, 최고급 화면과 공예를 추구하는 사용자, 스마트폰으로 영상 촬영하여 콘텐츠를 만드는 크리에이터.

\textbf{대표 기종}: iPhone 17 Pro/Pro Max, 삼성 S26 Ultra, 화웨이 Pura 80 Ultra, 샤오미 17 Ultra, vivo X300 Pro.

\textbf{구매 가이드}: 이 가격대에서 각 브랜드의 차이가 점점 더 작아지고 있어, 사실 주로 어떤 시스템을 좋아하느냐에 달려 있음. 촬영 좋아하면 vivo 화웨이, 영상 촬영 좋아하면 iPhone, 화면과 SPen 좋아하면 삼성, 균형 원하면 샤오미 선택. 솔직히 이 가격대에서는,\textbf{제멋대로 할 수 있음, 마음에 드는 것을 선택하면 됨}.

\booksection{시나리오별 기기 선택——당신의 요구는 무엇인가?}

다음은 가장 흔한 “사용자 프로필”로, 해당되는 것을 고르면 됨:

\begin{center}
\begin{tabularx}{\linewidth}{|X|X|X|}
\hline
\textbf{사용 시나리오} & \textbf{중점적으로 볼 것} & \textbf{추천 방향} \tabularnewline
\hline
사진 전문가 & 대형 센서 메인 카메라(1인치 최적), 광학 손떨림 보정, 망원(3x 이상 광학 줌), 영상 칩 & vivo X 시리즈, 화웨이 Pura/Mate 시리즈, 샤오미 Ultra 시리즈 \tabularnewline
게임 플레이어 & 최고급 프로세서, 높은 터치 샘플링, 방열 시스템, 대용량 배터리+급속 충전 & iQOO 숫자 시리즈, ROG Phone, Red Magic, Redmi K 시리즈 \tabularnewline
배터리 수명 우선 & 5500mAh 이상 배터리, 에너지 효율 좋은 프로세서, LTPO 화면 & Honor Magic 시리즈, OnePlus, Redmi Note 대용량 배터리 버전 \tabularnewline
어르신 사용 & 간편 모드, 큰 글꼴, 큰 볼륨, 대용량 배터리, A/S 서비스 지점 많음 & 화웨이 Changxiang, Honor Play, OPPO A 시리즈, Redmi Note \tabularnewline
학생 기기 & 가성비 높음, 내구성, 대용량 배터리, 학부모 통제 기능 & Redmi Note/Turbo, realme GT, iQOO Neo \tabularnewline
Apple 생태계 사용자 & 이미 Mac/iPad/AirPods 보유, iPhone 구매가 유일한 해법 & iPhone 17 표준판(가성비) 또는 Pro 버전(완전체 경험) \tabularnewline
해외 여행 & eSIM 지원, 글로벌 주파수 대역 완전, Google 서비스 & iPhone, 삼성 S 시리즈, Pixel, OnePlus \tabularnewline
\hline
\end{tabularx}
\end{center}

\booksection{다섯 가지 쉽게 빠지는 함정——구매 전 대비하여 한 번 보기}

\begin{enumerate}
  \item \textbf{함정 1: 4카메라 5카메라 $\neq$ 촬영 좋음}——제조사는 파라미터에 렌즈 수를 쌓는 것을 좋아하며, 그 안에는 자주 200만 화소의 "메꾸기 렌즈"(접사+배경 흐림 효과) 두 개가 있으며, 실제로는 완전히 쓸모 없음. 스마트폰 촬영 능력을 볼 때,\textbf{메인 카메라+초광각+망원}이 세 가지만 보면 충분. 나머지 추가는 없는 것으로 간주.
  \item \textbf{함정 2: 1억 화소에 현혹되지 말 것}——화소 수가 높다고 화질이 좋은 것이 아님. 센서 크기, 조리개, 알고리즘 튜닝이 화소보다 훨씬 중요. 많은 플래그십이 5000만 화소의 대형 센서를 사용하여, 찍은 사진이 일부 1억 화소의 천 위안 기기보다 훨씬 좋음.
  \item \textbf{함정 3: 통신사 커스텀 기기 구매 금지}——"요금 저장 후 스마트폰 증정"이 보기에 매우 합산해 보이지만, 통신사 커스텀 기기는 시스템 업데이트가 느리고, 삭제할 수 없는 사전 설치 앱이 있을 수 있으며, 보통 통신사 잠금(해당 통신사 카드만 사용 가능)되어 있음. 돈은 아낀 것처럼 보이지만, 실제로 "잠금"당함.
  \item \textbf{함정 4: "곡면 화면이 매우 고급스럽다"——마케팅이지 진리가 아님}——곡면 화면은 한때 플래그십의 표지였지만, 이제 점점 더 많은 사람이 그 문제를 발견: 보호 필름 붙이기 어려움, 가장자리 오작동, 깨지면 수리비 비쌈. 2026년의 추세는 "미세 곡면" 또는 "완전 직사각형"으로 회귀. "곡면이 고급스럽다"라는 이유로 곡면 화면을 선택하지 말 것.
  \item \textbf{함정 5: 모든 "AI 스마트폰"이 진짜 유용한 것이 아님}——2026년 모든 제조사가 "AI 스마트폰"을 외치지만, 많은 것은 기존의 음성 어시스턴트 이름만 바꾼 것. 진짜 유용한 AI 기능은: 실시간 번역, AI 제거(촬영 시 사람 제거), AI 보정, 단말기 대규모 모델 추론, 현지화된 스마트 요약. 나머지 화려한 기능은 기본적으로 허세.
\end{enumerate}

\booksection{언제 사는 것이 가장 합산한가?}

\begin{itemize}
  \item \textbf{618(5월 말\textasciitilde{}6월 중순)}: 1년 중 스마트폰 가격 인하 폭이 가장 큰 시기. 많은 스마트폰 제조사가 매출 실적을 위해, 플래그십이 800-1,500위안 가격 인하 가능. 상반기 스마트폰 구매 최적의 창.
  \item \textbf{쌍11(10월 말\textasciitilde{}11월 중순)}: 하반기 대 프로모션, 강도와 618 거의 비슷. 이 기간에 많은 브랜드 차세대 플래그십이 막 발표되어, 구세대 플래그십 가격 인하 폭이 가장 큼.
  \item \textbf{신기종 발표 후 1개월}: 막 발표되었을 때 원가, 1개월 후 채널가(제3자 매장) 일반적으로 200-500 저렴.
  \item \textbf{춘절 전 2주부터 원소절까지}: 오프라인 실체 매장 재고 정리 "녠화제(년화물절)", 직원과 가격 협상 가능.
  \item \textbf{비추천 시기}: 신제품 발표회 당일 첫 구매——정말 첫 번째 사용자가 되고 싶지 않다면, 첫 구매 가격이 가장 비쌈.
\end{itemize}

\booksection{채널 선택——어디서 사는 것이 신뢰할 수 있는가?}

\begin{itemize}
  \item \textbf{브랜드 공식 몰/티몰 플래그십 스토어/JD 셀프 운영}: 가장 안전, 가격 우위 없지만 리퍼 및 가품 우려 없음, A/S 서비스 최고.
  \item \textbf{핀둬둬 "100억 보조금"}: 가격이 가장 낮음, 단 "100억 보조금"으로 표시된 매장을 선택해야 함. 표시되지 않은 것은 손대지 말 것. 상품 수령 시 개박스 영상 녹화 기억, 문제 있으면 반품 가능.
  \item \textbf{오프라인 실체 매장}: 직접 만지고 손맛 시험 가능. 주의: 가이드는 이익이 높은 기종(당신에게 가장 적합한 것이 아니라)을 추천할 것, 인터넷에서 2-3개 목표 기종을 먼저 확정한 후 가서 "나는 이 몇 개만 보고 싶다"고 말하는 것을 권장.
  \item \textbf{중고 플랫폼(시안위/주안쥐안)}: 저렴하지만 위험 있음. 중고를 사려면, 세 가지 기억: 구매 영수증 또는 주문 스크린샷 제공 요청, IMEI 번호 검수 요청,\textbf{플랫폼 거래만, 절대 사적 이체 금지}.
\end{itemize}

\begin{tcbpractice}{실습: 이 방법론으로 한 대의 스마트폰 선택}
\begin{enumerate}
  \item 종이 한 장을 꺼내, 두 가지 핵심 요구 적기(예: 배터리 수명 1순위, 촬영 2순위).
  \item 예산 상한 정하기(예: 3,500위안).
  \item JD 또는 티몰을 열고, 가격 필터로, 3-5개 후보 기종 동그라미.
  \item B스테이션에서 "XX 기종 한 달 사용 경험" 검색(첫날 리뷰가 아닌, 최소 한 달 사용 후의 진짜 감상), 두세 개 영상 보기.
  \item 채널 신뢰 확인 후 주문, 수령 시 개박스 영상 녹화.
\end{enumerate}

\end{tcbpractice}

\begin{tcbsummary}{본 챕터 요약}
\begin{itemize}
  \item 기기 선택 세 가지 도끼: 먼저 요구 예산 마지막에 기종 선택
  \item 핵심 파라미터 다섯 가지만 보기: 프로세서, 화면, 카메라, 배터리 충전, 시스템 업데이트 약속
  \item 1500-3500위안은 가성비 최고, 가장 추천하는 스윗 스팟 구간
  \item 작년 플래그십이 올해 중급보다 더 합산함
  \item 5대 함정: 다카메라 사기, 고화소 함정, 통신사 커스텀, 곡면 화면 미신, AI 허세
  \item 618과 쌍11은 1년 중 스마트폰을 사기에 가장 적합한 두 시간대
\end{itemize}

\end{tcbsummary}
